% Generated draft: requires mathematical and source review.
\documentclass[11pt]{article}
\usepackage[margin=25mm]{geometry}
\usepackage{fontspec}
\setmainfont{Latin Modern Roman}
\usepackage{amsmath,amssymb,mathtools,unicode-math}
\setmathfont{Latin Modern Math}
\usepackage{enumitem,needspace,xurl,hyperref}
\setlength{\parindent}{0pt}
\setlength{\parskip}{5pt}
\setlength{\emergencystretch}{3em}
\newcommand{\N}{\mathbb{N}}
\newcommand{\Z}{\mathbb{Z}}
\newcommand{\Q}{\mathbb{Q}}
\newcommand{\R}{\mathbb{R}}
\newcommand{\C}{\mathbb{C}}
\newcommand{\F}{\mathbb{F}}
\newcommand{\A}{\mathbb{A}}
\newcommand{\PP}{\mathbb{P}}
\newcommand{\E}{\mathbb{E}}
\newcommand{\eps}{\varepsilon}
\newcommand{\dd}{\,\mathrm{d}}
\begin{document}
\section*{UnsolvedMath ID: 30003499}
\textbf{Source problem code:} catalogue entry
\subsection{Definitions and mathematical statement}
The supplied problem asks for conditions ensuring the equalities
\[
A^\infty(X)=C^\infty(X),\qquad A^M(X)=C^M(X),
\]
for non-open subsets \(X\subseteq\R^d\). These are characterization questions, not assertions that the equalities hold for every such set. The source does not fully specify the four classes, so the following definitions explain a possible scalar-valued interpretation and identify the conventions still requiring confirmation.

A smooth arc in \(X\) is a map \(c:I\to\R^d\), where \(I\subseteq\R\) is an open interval, such that \(c\) is smooth and \(c(I)\subseteq X\). For real-valued functions, the usual smooth-arc testing class is
\[
\mathcal A^\infty(X)=\{f:X\to\R: f\circ c\in C^\infty(I)\text{ for every smooth arc }c:I\to X\}.
\]
One possible interpretation of smoothness on \(X\) is extension to an open neighborhood:
\[
\mathcal E^\infty(X)=\{f:X\to\R:\text{ there exist an open }U\supseteq X\text{ and }F\in C^\infty(U)\text{ with }F|_X=f\}.
\]
If the intended conventions are \(A^\infty(X)=\mathcal A^\infty(X)\) and \(C^\infty(X)=\mathcal E^\infty(X)\), the smooth question is to characterize those non-open \(X\) for which every function passing all smooth-arc tests admits such an extension. This interpretation is conditional: the supplied record does not specify the codomain or whether its Whitney-type class means extendible functions, functions admitting compatible jets, or jets themselves. A Whitney jet consists of functions \(f_\alpha:X\to\R\), indexed by derivative multi-indices \(\alpha\), subject to the appropriate Taylor remainder compatibility conditions; it contains more data than its zeroth component \(f_0\). Thus a comparison with a class of functions requires an explicit convention about passing from jets to their zeroth components.

For the ultradifferentiable question, \(M\) denotes an unspecified ultradifferentiable datum. Once an ambient class \(C^M\), a family \(\mathscr T^M(X)\) of testing arcs with their open parameter intervals, and a Whitney-type class on \(X\) have been fixed, the arc-testing construction has the form
\[
\mathcal A^M_{\mathscr T}(X)=\{f:X\to\R: f\circ c\in C^M(I)\text{ for every }(I,c)\in\mathscr T^M(X)\}.
\]
The intended task is to determine hypotheses on \(X\) and \(M\) under which this testing class agrees with the intended \(C^M(X)\). The record does not determine whether the testing arcs are smooth or ultradifferentiable, the growth estimates defining \(C^M\), the assumptions on \(M\), or the Roumieu versus Beurling convention. Consequently it does not determine a single precise ultradifferentiable equality without further source clarification.

The title refers to closed sets, whereas the statement explicitly allows non-open subsets without requiring closedness. Closedness must therefore remain an unresolved scope qualification, rather than an added hypothesis.
\subsection{Short English statement}
Determine when testing a function along arcs detects exactly the intended Whitney-type smoothness on a non-open subset of Euclidean space, both in the smooth and ultradifferentiable settings. The meanings of the Whitney-type classes, the ultradifferentiable datum \(M\), and the permitted testing arcs require clarification. The supplied record labels the problem partially solved; this draft does not independently verify that assessment.
\subsection{Sources}
\begin{itemize}
\item Supplied local UnsolvedMath registry record, permanent ID \(30003499\), source problem code \texttt{catalogue entry}, in the Oberwolfach Reports Open Problems collection. Title: \emph{Arc-Smooth and Whitney-Smooth Functions on Closed Sets}. No individual proposer is supplied; the proposed year is recorded as 2017.
\item Supplied source citation: \emph{O-Minimality and its Applications to Number Theory and Analysis} (2018), Oberwolfach Reports, DOI \texttt{10.4171/owr/2017/22}, pp.\ 1349--1420. \url{https://doi.org/10.4171/owr/2017/22}. The background describes Oberwolfach Reports as documenting mathematical workshops and their open-problem discussions. The literature-source label is ``Oberwolfach Reports 2018.'' The supplied 2018 designation and the 2017 proposed year and DOI component are retained without resolving their chronology.
\item Supplied literature reference: Rainer, \emph{arc-smooth functions on closed sets}. \url{https://arxiv.org/abs/1801.08335}.
\item Literature assessment dated exactly \texttt{2026-08-21}: the supplied curation labels the problem partially solved and reports results for important non-open domains, including closed fat sets and subanalytic settings, while describing arbitrary non-open subsets and the full ultradifferentiable framework as unresolved. This is supplied dated curation derived from the OWR record, not an independently verified assessment. Time-sensitive status and the cited mathematical claims require independent verification.
\end{itemize}
\end{document}
